\section{Fragility Score Construction}
\label{sec:fragility_score_construction}

The predictive models produce a continuous estimate of forward 63-day maximum drawdown for each eligible stock-quarter. To obtain an interpretable cross-sectional risk ranking, these predictions are transformed into a stock-level \emph{Fragility Score}. For stock $i$ in quarter $q$, the score is

\begin{equation}
    \mathrm{FS}_{iq}
    = 100\frac{r_{iq}-1}{N_q-1},
    \label{eq:fragility_score}
\end{equation}

\noindent
where $r_{iq}$ is the ascending rank of the stock's predicted drawdown among the $N_q$ eligible stocks in that quarter. Tied predictions receive their average rank; if a quarter contains only one eligible stock, its score is set to 50.

The score lies between 0 and 100. When the extreme predictions are unique, the lowest receives 0 and the highest receives 100; tied predictions receive the score implied by their average rank. A score of 90 therefore means that the stock is predicted to be more fragile than approximately 90\% of the eligible universe in the same quarter. Ranking is performed separately within each quarter and uses only the contemporaneous cross-section of model predictions.

This transformation preserves relative ordering rather than the predicted level of drawdown. It is therefore suited to downside-risk sorting and portfolio applications, while the accuracy of the underlying magnitude predictions is evaluated separately using the error metrics. The primary economic evaluation applies the score to the model selected using the validation sample.
